\subsection{Definities en eerste eigenschappen}

\begin{lemma}[2.1.3]
    Zij $V$ een $K$-vectorruimte. We noteren het neutraal element voor de optelling in $K$ met $0_K$ en het neutraal element voor de optelling in $V$ met $0_V$.
    \begin{enumerate}
        \item[(i)] Het neutraal element in $V$ voor de optelling in $V$ is uniek.
        \item[(ii)] Elke $v \in V$ heeft een uniek tegengestelde voor de optelling; we noteren het als $-v$.
        \item[(iii)] Voor alle $v \in V$ is $-(-v) = v$.
        \item[(iv)] Voor alle $\lambda \in K$ is $\lambda 0_V = 0_V$.
        \item[(v)] Voor alle $v \in V$ is $0_K v = 0_V$.
        \item[(vi)] Voor alle $v \in V$ is $(-1)v = -v$. Bovendien geldt, voor alle $v \in V$ en $\lambda \in K$, dat
        \[ (-\lambda)v = -(\lambda v) = \lambda(-v), \quad \lambda v = (-\lambda)(-v), \quad -(v + w) = -v + (-w). \]
        \item[(vii)] Als voor een $v \in V$ en $\lambda \in K$ geldt dat $\lambda v = 0$, dan is $\lambda = 0_K$ of $v = 0_V$.
    \end{enumerate}
\end{lemma}
\begin{proof}(i)
    Voor \(v \in V: \) \(v + 0 = v \), kies \(v=0'\), dan:
     \(0'+ 0 = 0'.\)

    Stel niet uniek:
    \(\exists 0' : v + 0' = v\) kies \(v = 0\), dan:
    \(0 + 0' = 0\)

    Er volgt: \(0 = 0' = 0_V\)
\end{proof}
\begin{proof}(ii)
    Stel niet uniek, dan \(\exists v', v'': v + v' = 0_V \quad \text{en} \quad v + v'' = 0_V\).
    
    We vertrekken uit:
    \begin{align*}
        v + v' 
        &= 0_V \\
        \implies
        v + v'+ v'' &= v'' \\
        (v+v'')+v' &= v'' \\
        v'&=v'' 
    \end{align*}
\end{proof}
\begin{proof}(iii)
        \(v + (-v) = 0_V\), neem \(v = (-v)\)

        \(\implies (-v) + (-(-v)) = 0_V\)
        \(\implies  v + (-v) + (-(-v)) = v + 0_V \)
        \(\implies -(-v) = v\)
\end{proof}
\begin{proof}(iv)
    \(
        0_V 
        = \lambda v + (- \lambda v) 
        = \lambda(v + (-v)) 
        = \lambda 0_V
    \)
\end{proof}
\begin{proof}(v)
    \(
        0_V 
        = \lambda v + (- \lambda v) 
        = (\lambda + (-\lambda))v 
        =  0_K v
    \)
\end{proof}
\begin{proof}(vi)
    \begin{align*}
        (1 + (-1)) &= 0_K \\
        (1 + (-1))v &= 0_V \\
        (1)v + (-1)v &= 0_V \\
        -v + v + (-1)v &= -v \\
        (-1)v = -v 
    \end{align*}
\end{proof}
\begin{proof}(vii)
    stel \(\lambda \neq 0_K: \exists \lambda^{-1}: \lambda \lambda^{-1} = \lambda^{-1} \lambda = 1\)

    Dan:
    \(\lambda v = 0 
    \implies \lambda^{-1}\lambda v = \lambda^{-1}0 = 0
        \implies v = 0
    \)
\end{proof}

\subsection{Deelruimten, lineaire combinaties, span}

\begin{lemma}[2.2.3]
    Zij $K$ een veld en $V$ een $K$-vectorruimte, en zij $\emptyset \neq W \subseteq V$ een deelverzameling van $V$. Dan is $W$ een deelruimte van $V$ als en slechts als
    \[ \lambda_1 w_1 + \lambda_2 w_2 \in W \]
    voor alle $w_1, w_2 \in W$ en voor alle $\lambda_1, \lambda_2 \in K$.
\end{lemma}
\begin{proof}

    \([\Leftarrow]\) Als aan \( \lambda_1 w_1 + \lambda_2 w_2 \in W \) voldaan is tonen we aan dat \(W\) een deelruimte is.

    We checken V1-V7:

    (V1) Triviaal.

    (V2) Neem \(\lambda_1 = \lambda_2 = 0 \implies 0 \in W \).

    (V3) Neem \(\lambda_1 = 1 \) en \( \lambda_2 = -1\).

    (V4) Triviaal voor optelling.

    (V5) Triviaal.

    (V6) Neem \(\lambda_1 = 1\) en \(\lambda_2 = 0\).

    (V7) 1) Neem \(w_1 = w_2\) en 2) neem \(\lambda_1 = \lambda_2\).

    \(W\) is inderdaad een deelruimte.

    \([\Rightarrow]\) Omgekeerd: Als \(W\) een deelruimte is van \(V\) dan voldoet 
    \( \lambda_1 w_1 + \lambda_2 w_2 \in W .\)

    Dit is per definitie van de uniciteit voor een vectorruimte.
\end{proof}

\begin{lemma}[2.2.9]
    Zij $V$ een $K$-vectorruimte, en zij $S \subseteq V$.
    \begin{enumerate}
        \item[(i)] De verzameling $\operatorname{span}(S)$ is een deelruimte van $V$.
        \item[(ii)] Als $S \subseteq T \subseteq V$, dan is $\operatorname{span}(S) \subseteq \operatorname{span}(T)$.
    \end{enumerate}
\end{lemma}
\begin{proof}(i)
    \(\operatorname{span}(S) = \{\lambda_1s_1+ \dots + \lambda_is_i + \dots |s_i \in S, \lambda_i \in K \}\)

    \([\Rightarrow]\) Stel \(S \subseteq V\), dan is \(\operatorname{span}(S) \neq \emptyset\) een deelruimte van \(V\) als \(\operatorname{span}(S)\) een vectorruimte is.

    Om aan te tonen dat \(\operatorname{span}(S)\) een vectorruimte is moet die voldoen aan \(\lambda_1 w_1 + \lambda_2 w_2 \in \operatorname{span}(S)\) met \(w_1,w_2 \in \operatorname{span}(S) \)
    en \(\lambda_1,\lambda_2 \in K\). 

    Met \(w_1 = \sum_{i} \nu_i s_i\) en \(w_2 = \sum_{j} \mu_j s_j\).

    Dan: 
    \(\lambda_1 w_1 + \lambda_2 w_2 = \lambda_1 \sum_{i} \nu_i s_i + \lambda_2 \sum_{j} \mu_j s_j = \sum_{k=i+j} \gamma_k s_k \in \operatorname{span}(S) \).    
\end{proof}
\begin{proof}(ii)
    \([\Rightarrow]\) Als \(S \subseteq T\) dan is \(\operatorname{span}(S) = \sum_{i} \lambda_i s_i \in S \)
    en \(\operatorname{span}(T) = \left[\sum_{i} \lambda_i s_i + \sum_{i} \mu_i t_i \right] \in T\) 
    waarbij \(t_i \notin S\) maar wel \(t_i \in T\). Bijgevolg is \(\operatorname{span}(S) \subseteq \operatorname{span}(T)\).
\end{proof}

\subsection{Lineaire (on)afhankelijkheid en voortbrengendheid}

\begin{lemma}[2.3.6]
    Zij $V$ een $K$-vectorruimte.
    \begin{enumerate}
        \item[(i)] Zij $\{v_1, \dots, v_n\} \subseteq V$, en zij $\lambda_1 v_1 + \dots + \lambda_n v_n = 0$ een lineaire combinatie met $\lambda_i \neq 0$ voor een zekere $i \in \{1, \dots, n\}$. Dan is $v_i$ een lineaire combinatie van de overige $n - 1$ elementen.
        \item[(ii)] Twee elementen in $V$ zijn lineair afhankelijk als en slechts als de ene een scalair veelvoud van de andere is.
    \end{enumerate}
\end{lemma}
\begin{proof}(i)
    Er volgt dat: \(-\lambda_i v_i = \lambda_1 v_1 + ... +\lambda_n v_n\)
    \(\implies v_i = -\sum_{j}^{n-1} \frac{\lambda_j}{\lambda_i} v_j  \).
\end{proof}
\begin{proof}(ii)
    \([\Rightarrow]\) Stel \(v_1,v_2 \in V\) zijn lineair afhankelijk.
    Dit wil zeggen dat \(\mu_1 v_1 + \mu_2 v_2 = 0\) met \(v_1, v_2 \neq 0\) en \(\mu_1, \mu_2 \in K\).

    Er volgt dat \(v_1 = -\frac{\mu_2}{\mu_1}v_2\).
    
    \([\Leftarrow]\) Omgekeerd \(v_1 = c v_2\) een veelvoud;
    dan kunnen we op dezelfde manier redeneren dat nu \(v_1 - c v_2 = 0\) waarbij beide voorfactoren niet gelijk zijn aan 0.
\end{proof}

\subsection{Basissen, dimensie}

\begin{stelling}[2.4.4]
    Een deelverzameling $\mathcal{B} \subseteq V$ is een basis voor $V$ als en slechts als elk element $v \in V$ op unieke manier te schrijven is als een lineaire combinatie van de elementen uit $\mathcal{B}$.
\end{stelling}
\begin{proof}
    \([\Rightarrow]\) Stel $\mathcal{B} \subseteq V$ een basis voor $V$, per definitie is \(\mathcal{B} \) voortbrengend voor \(V\).
    
    Uniciteit: stel \(v = \sum_{i}^{n} \mu_i b_i = \sum_{i}^{n} \nu_i b_i\) met minstens één \( \mu_i \neq 0 \neq \nu_i\) op twee verschillende manieren te schrijven.
    Dan geldt:
    \[
        \sum_{i}^{n} \mu_i b_i - \sum_{i}^{n} \nu_i b_i = 0
        \implies \sum_{i}^{n} (\mu_i-\nu_i)b_i = 0
    \]
    Wegens het lineair onafhankelijk zijn van de veronderstelde basis zijn alle \(\mu_i - \nu_i = 0 \implies \mu_i = \nu_i\).

    \([\Leftarrow]\) Stel omgekeerd elk element $v \in V$ is op unieke manier te schrijven als lineaire combinatie van de elementen uit $\mathcal{B}$.

    Stel het nul element is te schrijven als \(\sum_i^n \mu_i b_i = 0 \).
    We kunnen ook schrijven \(0 = \sum_i^n 0 b_i \).

    Wegens de veronderstelde uniciteit is elke \(\mu_i = 0\), dus zijn de elementen van \(\mathcal{B}\) lineair onafhankelijk.
    Het voortbrengendheid volgt direect uit het gestelde.
\end{proof}

\begin{lemma}[2.4.6 (Lemma van Steinitz)]
    Zij $V$ een eindig-dimensionale $K$-vectorruimte met $S$ een eindige voortbrengende verzameling, $|S| = n \in \mathbb{N}$. Dan heeft elke lineair onafhankelijke deelverzameling van $V$ hoogstens $n$ elementen.
\end{lemma}
\begin{proof}
    Gegeven is dat $S = \{s_1, \dots, s_n\}$ voortbrengend is voor $V$, dus $\operatorname{span}(S) = V$. Elk element in $V$ kan geschreven worden als een lineaire combinatie van elementen uit $S$.

    Stel dat we een deelverzameling $U = \{u_1, \dots , u_{n+1}\}$ hebben met $n+1$ elementen. We zullen aantonen dat deze altijd lineair afhankelijk is.
    Dit wil zeggen dat we een niet-triviale oplossing zoeken voor:
    \[
    \sum_{i=1}^{n+1} \gamma_i u_i = 0 \quad \text{waarbij niet alle } \gamma_i = 0.
    \]
    Omdat $S$ de ruimte voortbrengt en $U \subseteq V$, kan elke $u_i$ geschreven worden als een lineaire combinatie van $S$:
    \[
    u_i = \sum_{j=1}^{n} \mu_{ij} s_j \quad \text{voor } i \in \{1, \dots, n+1\} 
    \]
    We vullen dit in onze oorspronkelijke som in:
    \[
        \sum_{i=1}^{n+1} \gamma_i u_i 
        = \sum_{i=1}^{n+1} \gamma_i \left( \sum_{j=1}^{n} \mu_{ij} s_j \right) 
        = \sum_{i=1}^{n+1} \sum_{j=1}^{n} \gamma_i \mu_{ij} s_j 
        = \sum_{j=1}^{n} \left( \sum_{i=1}^{n+1} \gamma_i \mu_{ij} \right) s_j = 0
    \]
    Om te zorgen dat deze totale lineaire combinatie nul wordt,
    is het voldoende als de coëfficiënt voor elke $s_j$ gelijk is aan nul.
    We zoeken dus oplossingen $\gamma_i$ zodanig dat:
    \[ \sum_{i=1}^{n+1} \gamma_i \mu_{ij} = 0 \quad \text{voor elke } j \in \{1, \dots, n\} \]
    Dit levert een homogeen stelsel op van $n$ vergelijkingen ($\forall j$) in $n+1$ onbekenden (de $\gamma_i$'s). 
    Omdat dit stelsel meer onbekenden dan vergelijkingen heeft,
    heeft het gegarandeerd minstens één vrije variabele.
    Er bestaat dus een niet-triviale oplossing (minstens één $\gamma_i \neq 0$).

    Omdat we niet-nul scalairen $\gamma_i$ hebben gevonden waarvoor $\sum_{i=1}^{n+1} \gamma_i u_i = 0$,
    is de verzameling $U$ per definitie lineair afhankelijk.
    Bijgevolg kan een lineair onafhankelijke verzameling in $V$ nooit meer dan $n$ elementen bevatten.
\end{proof}

\begin{stelling}[2.4.7]
    Zij $V$ een eindig-dimensionale $K$-vectorruimte. Zij $T$ een lineair onafhankelijke deelverzameling van $V$ en $S$ een voortbrengende verzameling die $T$ bevat. Dan is er een basis $\mathcal{B}$ voor $V$ zodat $T \subseteq \mathcal{B} \subseteq S$.
\end{stelling}
\begin{proof}
    Als \(T\) niet voortbrengend is, dan kunnen we deze aanvullen met een lineair onafhankelijk element uit \(S\)
    tot de nieuwe verzameling \(\mathcal{B}\) wel voortbrengend is en dus een basis vormt waarbij \(T \subseteq \mathcal{B}\).
    
    Stel dat dit niet zo is, dan zou \(T \cup s \in S\) lineair afhankelijk zijn. 
    Dan zouden we \(\sum \nu_i t_i + \mu_j s = 0\) kunnen vinden waarbij niet alle \(\nu_i, \mu_j = 0\),
    en \(s = -\sum \frac{\nu_i}{\mu_j} t_i \in \operatorname{Span}(T)\). 
    Omdat \(s \in S\) willekeurig gekozen werd betekend dit dat \(S \subseteq\operatorname{Span}(T)\).

    \(\implies S = T\), dus \(T\) voortbrengend. Dit is in strijd met het gestelde. 

    Is \(S\) niet lineair onafhankelijk,
    dan vinden we een combinatie \(s_i = -\sum_{j} \frac{\mu_j}{\mu_i}s_j\).
    Waarbij deze \(s_i\) verwijderd kan worden zonder de voortbrengendheid te schenden. 
    Deze actie kan opnieuw gedaan worden zolang \(S\) lineair afhankelijk is.
\end{proof}

\begin{gevolg}[2.4.8]
    \begin{enumerate}
        \item[(i)] Zij $V$ een eindig-dimensionale $K$-vectorruimte, en zij $T \subseteq V$ een lineair onafhankelijke verzameling. Dan bestaat er een basis $\mathcal{B}$ voor $V$ zodat $T \subseteq \mathcal{B}$.
        \item[(ii)] Zij $V$ een eindig-dimensionale $K$-vectorruimte, en zij $S \subseteq V$ een voortbrengende verzameling. Dan bestaat er een basis $\mathcal{B}$ voor $V$ zodat $\mathcal{B} \subseteq S$.
    \end{enumerate}
\end{gevolg}
\begin{proof}(i)
    Volgt uit stelling 2.4.7 door het systematisch toevoegen van lin. onafh. elementen uit \(V\).
\end{proof}
\begin{proof}(ii)
    Volgt direct uit stelling 2.4.7 door het reduceren vd lineair afhankelijke elementen uit \(S\).
\end{proof}

\begin{gevolg}[2.4.9]
    Zij $V \neq \{0\}$ een eindig-dimensionale $K$-vectorruimte.
    \begin{enumerate}
        \item[(i)] Er is een basis in $V$.
        \item[(ii)] Elke basis in $V$ is eindig, en alle basissen hebben evenveel elementen.
    \end{enumerate}
\end{gevolg}
\begin{proof}(i)
    Omdat \(V\) eindig-dimensionaal is, heeft \(V\) per definitie een eindige voortbrengende verzameling \(S\).
    Als \(S\) lineair afhankelijk is, pas gevolg 2.4.8 (ii) systematisch toe tot \(S\) lineair onafhankelijk is.
    Dit geeft direct een basis \(\mathcal{B} \subseteq S\).
\end{proof}
\begin{proof}(ii)
    Stel basis \(\mathcal{B}\) en een andere basis \(\mathcal{C}\).
    Wegens voortbrengendheid van de basissen volgt uit lemma 2.4.6 dat:
      \[
      |\mathcal{C}| \geq |\mathcal{B}| \quad\text{en}\quad |\mathcal{B}| \geq |\mathcal{C}|
      \implies |\mathcal{C}| = |\mathcal{B}|
      \]
\end{proof}

\begin{stelling}[2.4.12]
    Zij $V$ een $n$-dimensionale $K$-vectorruimte, en zij $S \subseteq V$ een verzameling met precies $n$ elementen.
    \begin{enumerate}
        \item[(i)] Als $S$ lineair onafhankelijk is, dan is $S$ een basis.
        \item[(ii)] Als $S$ voortbrengend is, dan is $S$ een basis.
    \end{enumerate}
\end{stelling}
\begin{proof}(i)
    Uit gevolg 2.4.8 (i) volgt dat er een basis \(\mathcal{B}\) bestaat zodat \(S \subseteq \mathcal{B}\).

    \(\operatorname{Dim}(V) = n \implies |\mathcal{B}| = n\) en \(|S| = n \implies |\mathcal{B}| = |S| \implies \mathcal{B} = S\).
\end{proof}
\begin{proof}(ii)
    Uit gevolg 2.4.8 (ii) volgt dat er een basis \(\mathcal{B}\) bestaat zodat \( \mathcal{B} \subseteq S\).

    \(\operatorname{Dim}(V) = n \implies |\mathcal{B}| = n \) en \(|S| = n \implies |\mathcal{B}| = |S| \implies \mathcal{B} = S\).
\end{proof}

\subsection{Som en directe som van vectorruimten}

\begin{lemma}[2.5.2]
    Zij $V$ een $K$-vectorruimte en $W_1, W_2, W$ drie deelruimten van $V$. Dan is $W = W_1 \oplus W_2$ als en slechts als $W = W_1 + W_2$ en $W_1 \cap W_2 = \{0\}$.
\end{lemma}
\begin{proof}
    \([\Rightarrow]\) Stel \(W =W_1 \oplus W_2\):
    Per definitie is \(W = W_1 + W_2\).

    Stel \(w \in W_1 \cap W_2\), dan:
    \(0 = w + (-w)\) en \(0 = 0 + 0\) want \(0 \in W_1 \cap W_2 \).
    Door de uniciteit moet \(w=0\),
    Waardoor  \(W_1 \cap W_2 = \{0\} \).

    \([\Leftarrow]\) Stel omgekeerd dat $W = W_1 + W_2$ en $W_1 \cap W_2 = \{0\}$:

    Stel dat \(w = w_1 + w_2 = w_1' + w_2' \implies w_1 - w_1' = w_2' - w_2 \in W_1 \cap W_2\).

    Doordat \( W_1 \cap W_2 = 0\) is \(w_1 - w_1' =0 \implies w_1 = w_1'\) en \( w_2' - w_2= 0 \implies w_2 = w_2'\).
    
    Dus \(w = w_1 + w_2\) is uniek.
    Wegens \(W = W_1 + W_2\) is \(W = W_1 \oplus W_2\)
\end{proof}

\begin{stelling}[2.5.4]
    Zij $V$ een $K$-vectorruimte en $W_1, W_2, W$ drie deelruimten van $V$ zodat $W = W_1 \oplus W_2$.
    \begin{enumerate}
        \item[(i)] Als $\mathcal{B}_1$ een basis is voor $W_1$ en $\mathcal{B}_2$ een basis is voor $W_2$, dan is $\mathcal{B}_1 \cup \mathcal{B}_2$ een basis voor $W$.
        \item[(ii)] $\dim W = \dim W_1 + \dim W_2$.
    \end{enumerate}
\end{stelling}
\begin{proof}(i)
    Stel \(\mathcal{B} = \{b_1, \dots, b_n\}\) basis voor \(W_1\) en \(\mathcal{C} = \{c_1, \dots, c_m\}\) basis voor \(W_2\).
    dan is \(\mathcal{B} \cup \mathcal{C} = \{b_1, \dots, b_n,c_1, \dots, c_m\}\).

    Nu moeten we bewijzen dat dit een basis is voor \(W = W_1 \oplus W_2\).

    \[\sum_{i}^{n} \mu_i b_i + \sum_{j}^{m} \nu_j c_j  = 0\]

    Doordat \(W_1 \cap W_2 = \{0\}\) is \(\sum_{i}^{n} \mu_i b_i = -\sum_{j}^{m} \nu_j c_j \implies \mu_i = \nu_j = 0 , \,\forall i,j \in \{0,\dots,\max(n,m)\} \).

    Door dat dit geldt voor de basissen apart, geldt de lineair onafhankelijkheid ook voor de som van de basissen.
    De voortbrengendheid volgt uit dat \(\forall w_1 \in W_1: w_{1} = \sum_{i}^{n} \mu_i b_i \) en
     \(\forall w_2 \in W_2: w_{2} = \sum_{j}^{m} \mu_j b_j \).

    Uit \(W = W_1 + W_2\) volgt \(\forall w \in W: w = \sum_{i}^{n} \mu_i b_i + \sum_{j}^{m} \nu_j c_j\).
\end{proof}
\begin{proof}(ii)
    \(\dim(W)= |\mathcal{B} \cup \mathcal{C}| \) het aantal elementen van de basis voor \(W\), doordat \(\mathcal{B} \cap \mathcal{C} = \{0\}\).  
    
    \(\dim(W)=  \dim(\mathcal{B} \cup \mathcal{C}) = \dim(\mathcal{B}) + \dim(\mathcal{C}) = \dim W_1 + \dim W_2. \)
\end{proof}

\begin{stelling}[2.5.5]
    Zij $V$ een $K$-vectorruimte met basis $\mathcal{B}$,
    en veronderstel dat $\mathcal{B}$ de disjuncte unie is van $\mathcal{B}_1$ en $\mathcal{B}_2$,
    i.e.\ $\mathcal{B} = \mathcal{B}_1 \cup \mathcal{B}_2$ en $\mathcal{B}_1 \cap \mathcal{B}_2 = \emptyset$.
    Dan is $V = \operatorname{span}(\mathcal{B}_1) \oplus \operatorname{span}(\mathcal{B}_2)$.
\end{stelling}
\begin{proof}
    \(V = \operatorname{span}(\mathcal{B}), V_1 = \operatorname{span}(\mathcal{B}_1) \) en \(V_2 = \operatorname{span}(\mathcal{B}_2)\).

    We moeten aantonen dat $V_1 \cap V_2 = \{0\}$.
    Stel dat $z \in V_1 \cap V_2$.

    Omdat $z \in V_1$, is $z = \sum_{v \in \mathcal{B}_1} \lambda_v v$ .
    Omdat $z \in V_2$, is $z = \sum_{w \in \mathcal{B}_2} \mu_w w$.
    \[\implies \sum_{v \in \mathcal{B}_1} \lambda_v v = \sum_{w \in \mathcal{B}_2} \mu_w w \implies \sum_{v \in \mathcal{B}_1} \lambda_v v - \sum_{w \in \mathcal{B}_2} \mu_w w = 0 \]
    Omdat $\mathcal{B} = \mathcal{B}_1 \cup \mathcal{B}_2$ een basis is, zijn de elementen in $\mathcal{B}$ lineair onafhankelijk.
     Aangezien bovendien $\mathcal{B}_1 \cap \mathcal{B}_2 = \emptyset$,
    moeten alle coëfficiënten $\lambda_v$ en $\mu_w$ noodzakelijk gelijk zijn aan 0.
    Bijgevolg is $z = 0$, wat bewijst dat $V_1 \cap V_2 = \{0\}$.

    \(\mathcal{B} = \mathcal{B}_1 + \mathcal{B}_2 \implies
     \operatorname{span} (\mathcal{B})= \operatorname{span} (\mathcal{B}_1) + \operatorname{span}(\mathcal{B}_2)
     \implies V = V_1 + V_2 \)

     Bijgevolg is \(V = V_1 \oplus V_2 = \operatorname{span} (\mathcal{B}_1) \oplus \operatorname{span}(\mathcal{B}_2) \).
\end{proof}

\begin{stelling}[2.5.6 (Dimensiestelling voor deelruimten)]
    Zij $V$ een eindig-dimensionale $K$-vectorruimte en $W_1, W_2$ twee deelruimten van $V$. Dan is
    \[ \dim(W_1 + W_2) = \dim W_1 + \dim W_2 - \dim(W_1 \cap W_2). \]
\end{stelling}
\begin{proof}
    Stel $\dim(W_1 \cap W_2) = k$. Kies een basis voor de doorsnede $W_1 \cap W_2$:
    $$ \mathcal{B}_0 = \{u_1, \dots, u_k\} $$

    Via gevolg 2.4.8 kunnen we deze basis uitbreiden naar:
    \begin{itemize}
        \item Een basis voor $W_1$: $\mathcal{B}_1 = \{u_1, \dots, u_k, v_1, \dots, v_r\}$, zodat $\dim W_1 = k + r$.
        \item Een basis voor $W_2$: $\mathcal{B}_2 = \{u_1, \dots, u_k, w_1, \dots, w_s\}$, zodat $\dim W_2 = k + s$.
    \end{itemize}

    We tonen nu aan dat de verzameling $\mathcal{B} = \{u_1, \dots, u_k, v_1, \dots, v_r, w_1, \dots, w_s\}$ een basis vormt voor de somruimte $W_1 + W_2$.

    $\mathcal{B}$ is voortbrengend voor $W_1 + W_2$:
    Zij $x \in W_1 + W_2$. 
    
    Per definitie $\exists x_1 \in W_1$ en $x_2 \in W_2: x = x_1 + x_2$.
    Omdat $x_1 \in \operatorname{span}(\mathcal{B}_1)$ en $x_2 \in \operatorname{span}(\mathcal{B}_2)$:
    $$ x_1 = \sum_{i=1}^k \alpha_i u_i + \sum_{j=1}^r \beta_j v_j \quad \text{en} \quad x_2 = \sum_{i=1}^k \alpha'_i u_i + \sum_{m=1}^s \gamma_m w_m $$
    Dan geldt:
    $$ x = x_1 + x_2 = \sum_{i=1}^k (\alpha_i + \alpha'_i) u_i + \sum_{j=1}^r \beta_j v_j + \sum_{m=1}^s \gamma_m w_m $$
    Hiermee is $x$ geschreven als een lineaire combinatie van elementen uit $\mathcal{B}$, dus $\operatorname{span}(\mathcal{B}) = W_1 + W_2$.

    $\mathcal{B}$ is lineair onafhankelijk:
    Stel dat er scalairen $\alpha_i, \beta_j, \gamma_m \in K$ bestaan zodanig dat:
    $$ \sum_{i=1}^k \alpha_i u_i + \sum_{j=1}^r \beta_j v_j + \sum_{m=1}^s \gamma_m w_m = 0 
     \implies \sum_{i=1}^k \alpha_i u_i + \sum_{j=1}^r \beta_j v_j = -\sum_{m=1}^s \gamma_m w_m $$
    Het linkerlid is een element van $W_1$ (combinatie uit $\mathcal{B}_1$). Het rechterlid is een element van $W_2$ (combinatie uit $\mathcal{B}_2$). 
    Dit betekent dat de vector $y = -\sum_{m=1}^s \gamma_m w_m$ tot de doorsnede $W_1 \cap W_2$ behoort.
    Omdat $y \in W_1 \cap W_2$:
    $$ y = -\sum_{m=1}^s \gamma_m w_m = \sum_{i=1}^k \delta_i u_i 
    \implies \sum_{i=1}^k \delta_i u_i + \sum_{m=1}^s \gamma_m w_m = 0 $$
    Aangezien $\{u_1, \dots, u_k, w_1, \dots, w_s\} = \mathcal{B}_2$ een basis is voor $W_2$, 
    zijn deze vectoren lineair onafhankelijk, $\implies \gamma_m = 0$ voor alle $m \in \{1, \dots, s\}$.
    Als we $\gamma_m = 0$ invullen in onze eerste combinatie, houden we over:
    $$ \sum_{i=1}^k \alpha_i u_i + \sum_{j=1}^r \beta_j v_j = 0 $$
    Omdat $\mathcal{B}_1 = \{u_1, \dots, u_k, v_1, \dots, v_r\}$ een basis is voor $W_1$,
    zijn ook deze vectoren lineair onafhankelijk,  $\implies \alpha_i = 0$ (voor alle $i$) en $\beta_j = 0$ (voor alle $j$).

    Alle coëfficiënten zijn noodzakelijk nul, wat bewijst dat $\mathcal{B}$ lineair onafhankelijk is.

    \(\implies \mathcal{B}\) vormt een basis voor $W_1 + W_2$. Het aantal elementen in $\mathcal{B}$ bepaalt de dimensie:
    $$ \dim(W_1 + W_2) = |\mathcal{B}| = k + r + s $$
    Ook is:
    $$ \dim W_1 + \dim W_2 - \dim(W_1 \cap W_2) = (k + r) + (k + s) - k = k + r + s $$
    Hieruit volgt dat:
    $$ \dim(W_1 + W_2) = \dim W_1 + \dim W_2 - \dim(W_1 \cap W_2) $$
\end{proof}

\begin{lemma}[2.5.8]
    Zij $W, U$ deelruimten van een $K$-vectorruimte $V$. Als $V = W + U$ dan bestaat er een deelruimte $Y \subseteq U$ zodat $V = W \oplus Y$. 
    In het bijzonder heeft elke deelruimte $W \subseteq V$ een complement in $V$.
\end{lemma}
\begin{proof}
    Als \(W \cap U = \{0\}\) dan is het gevraagde bewezen met \(Y = U\).

    Stel nu \(W \cap U \neq \{0\}\):
    \(\exists\) een basis \(\mathcal{B}_0 = (v_1,\dots,v_k)\) voor \(W \cap U\).

    Neem \(\mathcal{B}_1 = \mathcal{B}_0 \cup (w_1,\dots,w_{n-k}) \) een basis voor \(W\), \(\mathcal{B}_2 = \mathcal{B}_0 \cup (z_1,\dots,z_{m-k}) \) een basis voor \(U\) waarbij deze basissen de uitgebreide basissen zijn van \(\mathcal{B}_0\).

    Definieer nu de deelruimte $Y := \operatorname{span}(z_1, \dots, z_{m-k})$. 
    Omdat alle $z_i \in U$, geldt dat $Y \subseteq U$. We tonen aan dat $V = W \oplus Y$.

    Kies een willekeurige vector $x \in W \cap Y$.
    \begin{itemize}
        \item Omdat $x \in W$:
        \(x = \sum_{i=1}^{k} \mu_i v_i + \sum_{j=1}^{n-k} \nu_j w_j \)
        \item Omdat $x \in Y$:
        \( x = \sum_{l=1}^{m-k} \lambda_l z_l \)
    \end{itemize}
    Gelijkstellen geeft:
    \[ \sum_{i=1}^{k} \mu_i v_i + \sum_{j=1}^{n-k} \nu_j w_j = \sum_{l=1}^{m-k} \lambda_l z_l \implies \sum_{i=1}^{k} \mu_i v_i + \sum_{j=1}^{n-k} \nu_j w_j - \sum_{l=1}^{m-k} \lambda_l z_l = 0 \]
    Wegens (Stelling 2.5.6) is:
    \[ \mathcal{B} = \mathcal{B}_1 \cup \{z_1, \dots, z_{m-k}\} = \{v_1, \dots, v_k, w_1, \dots, w_{n-k}, z_1, \dots, z_{m-k}\} \]
    een basis voor $W + U = V$. 
    \( \implies \mu_i = 0, \quad \nu_j = 0, \quad \lambda_l = 0 \).
    $\implies x = 0$, wat bewijst dat $W \cap Y = \{0\}$.

    \(\implies V = \operatorname{span}(\mathcal{B}) = \operatorname{span}(\mathcal{B}_1) + \operatorname{span}(z_1, \dots, z_{m-k}) = W + Y \).
    
    Bijgevolg is \(V = W \oplus Y\).
\end{proof}

\begin{lemma}[2.5.11]
    Zij $W_1, \dots, W_m$ vectorruimten over $K$. De verzameling $\oplus_{i=1}^{m} W_i$ met optelling en scalaire vermenigvuldiging zoals gedefinieerd in Definitie 2.5.10, is een $K$-vectorruimte.
\end{lemma}
\begin{proof}
    Definieer op $\bigoplus_{i=1}^{m} W_i$ de volgende optelling en vermenigvuldiging met scalairen,
    \[ (a_1, \dots, a_m) + (b_1, \dots, b_m) = (a_1 + b_1, \dots, a_m + b_m) \]
    en
    \[ \lambda \cdot (a_1, \dots, a_m) = (\lambda a_1, \dots, \lambda a_m) \]
    voor alle $a_i, b_i \in W_i$ en $\lambda \in K$.

    Om te controlleren of deze ruimte een vectorruimte is, moeten we V1-V7 controlleren:

    (V1) Triviaal voor de optelling

    (V2) \(0_V = (0_1, \dots, 0_m)\).

    (V3) Neem \((b_1, \dots, b_m) = (-a_1, \dots, -a_m)\).

    (V4) Ja, optelling is commutatief.

    (V5) Ja, voor scalaire vermenigvuldiging geldt \((\lambda\mu)\cdot (a_1, \dots, a_m) = \lambda(\mu\cdot (a_1, \dots, a_m))\).

    (V6) Triviaal met \(\lambda = 1\).

    (V7) \((\lambda + \mu )\cdot (a_1, \dots, a_m) 
    =  ((\lambda + \mu )a_1, \dots, (\lambda + \mu )a_m)
    = (\lambda a_1 + \mu a_1, \dots, \lambda a_m + \mu a_m)
    = \lambda \cdot (a_1, \dots, a_m) +  \mu \cdot (a_1, \dots, a_m)\)

    en 
    
    \(
    \mu \cdot((a_1, \dots, a_m) + (b_1, \dots, b_m))
    = \mu \cdot (a_1 + b_1, \dots, a_m + b_m)
    = (\mu(a_1 + b_1), \dots, \mu(a_m + b_m))
    = (\mu a_1 + \mu b_1, \dots, \mu a_m + \mu b_m)
    = \mu \cdot (a_1, \dots, a_m) + \mu \cdot (b_1, \dots, b_m)
    \).
\end{proof}

\begin{lemma}[2.5.13]
    Zij $W_1, \dots, W_m$ vectorruimten over $K$, en stel $V = \oplus_{i=1}^{m} W_i$ de uitwendige directe som van $W_1, \dots, W_m$. Dan is $V$ de inwendige directe som van de vectorruimten
    \[ W_i' = \{(0, \dots, a_i, \dots, 0) \mid a_i \in W_i\}, \]
    i.e.\ $V = W_1' \oplus \dots \oplus W_m'$.
\end{lemma}
\begin{proof}
    Volgt uit \(\forall (a_1, \dots, a_m) \in \oplus_{i=1}^{m} W_i:\)

    \((a_1, a_2, \dots, a_m) = (a_1, 0, \dots, 0) + (0, a_2, \dots, 0) + \dots + (0, 0, \dots, a_m)\)

    de unieke manier is om \((a_1, a_2, \dots, a_m) \) te onbinden als som van elementen uit \(W_1', \dots, W_m'\).
\end{proof}
